\documentclass[11pt]{article}
\usepackage[a4paper,margin=2.2cm]{geometry}
\usepackage{amsmath,amssymb,amsfonts,bm}
\usepackage{physics}
\usepackage{hyperref}
\usepackage{enumitem}

\title{Notes: Single-site multiplets and 2nd-order perturbation for $f$-electron exchange (RS scheme)}
\author{}
\date{}

\begin{document}
\maketitle

\section{Scope and working assumptions}
We consider localized $f^n$ shells (single ion) and derive an effective two-site interaction via 2nd-order perturbation in inter-site hopping. The working hierarchy (typical $4f$):
\begin{equation}
H_{\text{int}} \gg H_{\text{SOC}} \gg H_{\text{CEF}},
\end{equation}
and we adopt a Russell--Saunders (LS) scheme with \emph{no cross-$LS$ mixing} under SOC at the current stage.

\section{Hamiltonians}
Single-ion Hamiltonian:
\begin{equation}
H_{\text{ion}} = H_{\text{int}} + H_{\text{SOC}} + H_{\text{CEF}}.
\end{equation}
Inter-site hopping (two sites $i,j$):
\begin{equation}
H_t=\sum_{p,q} t_{pq}\, f^\dagger_{i,p} f_{j,q} + \text{h.c.},
\qquad p,q\in\{0,\dots,13\}.
\end{equation}
We target the 2nd-order effective Hamiltonian in the low-energy subspace (ground manifold on each site):
\begin{equation}
H_{\text{eff}}^{(2)} = P H_t \frac{1}{E_0 - H_0} H_t P,
\end{equation}
with $H_0=\sum_\ell H_{\text{ion}}(\ell)$ and $P$ the projector to the low-energy (ground) subspace.

\section{Bases: LSMS and LSJM}
\subsection{LSMS (eigenstates of $H_{\text{int}}$)}
We denote LS-coupled multiplets (with possible multiplicity $\alpha$ for the same $LS$):
\begin{equation}
\ket{\alpha\,L M_L;\, S M_S}.
\end{equation}
Key point: when $\alpha>1$, the explicit eigenvectors within the $(L,S)$ subspace depend on the detailed $H_{\text{int}}$ (i.e.\ Slater--Condon/Racah parameters), so $\alpha$ is fixed by diagonalizing $H_{\text{int}}$ restricted to that $(L,S)$ sector.

\subsection{LSJM (SOC eigenstates within a fixed $\alpha LS$)}
In the current RS/no-$LS$-mixing setting, SOC only splits different $J$ \emph{within the same} $(\alpha,L,S)$:
\begin{equation}
\ket{\alpha\,L S;\, J M} = \sum_{M_L,M_S} C^{JM}_{L M_L,\, S M_S}\,\ket{\alpha\,L M_L;\, S M_S}.
\end{equation}
If $H_{\text{SOC}}=\lambda\,\bm L\cdot \bm S$, then within a fixed $(\alpha,L,S)$ the SOC energy is
\begin{equation}
E_{\text{SOC}}(\alpha LSJ)=\lambda\,\frac{J(J+1)-L(L+1)-S(S+1)}{2}.
\end{equation}
(If a paper uses $H_{\text{SOC}}=\frac{\lambda}{2}\bm L\cdot\bm S$, the factor becomes $1/4$ accordingly.)

\section{Coulomb parameters: ratio vs scale}
For $f$ shells, the interacting part is parameterized by $F^0,F^2,F^4,F^6$.
\begin{itemize}[leftmargin=2em]
\item The \textbf{wavefunctions} (i.e.\ the $\alpha$-resolved eigenvectors in LSMS, and hence LSJM) depend on the \textbf{ratios}:
\[
r_4=\frac{F^4}{F^2},\qquad r_6=\frac{F^6}{F^2}.
\]
Scaling $(F^2,F^4,F^6)\to s(F^2,F^4,F^6)$ rescales Coulomb energies but does not change eigenvectors (up to possible rotations inside exact degeneracies).
\item $F^0$ mainly provides an overall shift within a fixed $n$ sector and does not affect eigenvectors (it matters for absolute energies and energy \emph{differences across} $n$ sectors in perturbation denominators).
\end{itemize}

Thus it is useful to separate:
\begin{equation}
\text{shape}:(r_4,r_6) \quad \text{vs}\quad \text{scale}:F^2\ (\text{and }F^0).
\end{equation}

\section{CEF and the ground manifold}
After SOC splitting, CEF is applied in the (typically lowest) LSJ manifold:
\begin{equation}
H_{\text{CEF}} = \sum_{k,q} B_k^q\, O_k^q \quad (\text{Stevens operators}),
\end{equation}
diagonalized within the chosen LSJ subspace to obtain the actual ground manifold.
We denote the resulting ground subspace for $f^n$ as
\begin{equation}
\{\ket{g_a}\}_{a=1}^d,\qquad d=\text{degeneracy determined by the spectrum (not user-chosen)}.
\end{equation}
Each $\ket{g_a}$ is stored as a linear combination in the \emph{single-site Fock basis}.

\section{Fock basis and fermionic signs}
\subsection{Single-site Fock basis}
Let the 14 spin-orbitals be indexed by $p=0,\dots,13$. A determinant (occupation configuration) is a bitset $\ket{I}$.
Creation operator action:
\begin{equation}
f^\dagger_p \ket{I}=
\begin{cases}
0,& n_p(I)=1,\\
(-1)^{N_{<p}(I)} \ket{I\cup\{p\}},& n_p(I)=0,
\end{cases}
\end{equation}
where $N_{<p}(I)$ is the number of occupied orbitals with index $<p$.

Define sparse matrices (single-site):
\begin{equation}
D^{(n)}_p \equiv [f^\dagger_p]: \mathcal{B}_n \to \mathcal{B}_{n+1},
\qquad
C^{(n)}_p \equiv [f_p]: \mathcal{B}_n \to \mathcal{B}_{n-1},
\end{equation}
with the hermiticity relation in an orthonormal determinant basis:
\begin{equation}
C^{(n)}_p = \left(D^{(n-1)}_p\right)^\dagger.
\end{equation}

\subsection{Two-site ordering and parity string (formal point)}
Fix a global two-site orbital ordering:
\[
(i,0\ldots13)\prec(j,0\ldots13).
\]
Then operators on the $j$ site carry the parity string of site $i$:
\begin{equation}
f_{j,q} = P_i \otimes f_q,\qquad P_i=(-1)^{N_i}.
\end{equation}
In the 2nd-order superexchange processes with fixed $n$ on each site, the relevant parity factors appear twice (go and return) and cancel. Practically, we enforce a fixed ordering ID and keep contractions consistent.

\section{Intermediate states and projected transition operators}
For superexchange from $f^n\otimes f^n$, intermediate sectors are
\[
f^{n+1}\otimes f^{n-1}\quad \text{and}\quad f^{n-1}\otimes f^{n+1}.
\]
We neglect CEF in intermediate states (approximation):
\[
H_{\text{mid}} \approx H_{\text{int}} + H_{\text{SOC}},
\]
and use the LSJM eigenstates in each $f^{n\pm1}$ sector as the intermediate basis.

Let $V_{n+1}$ be the matrix whose columns are intermediate eigenstates in the $f^{n+1}$ sector expressed in the $f^{n+1}$ Fock basis:
\[
V_{n+1}\in \mathbb{C}^{\dim_{n+1}\times N_{n+1}},\qquad
V_{n-1}\in \mathbb{C}^{\dim_{n-1}\times N_{n-1}}.
\]
Define \textbf{projected transition operators} (single-site) used in SOPT:
\begin{align}
T^{(+)}_p &\equiv V_{n+1}^\dagger\, D_p^{(n)} \in \mathbb{C}^{N_{n+1}\times \dim_n},\\
T^{(-)}_q &\equiv V_{n-1}^\dagger\, C_q^{(n)} = V_{n-1}^\dagger \left(D_q^{(n-1)}\right)^\dagger
\in \mathbb{C}^{N_{n-1}\times \dim_n}.
\end{align}
This realizes the philosophy: \emph{hopping acts on the Fock basis; intermediate eigenstates enter only via linear-algebra projections}.

Ground manifold vectors are stored as
\[
G \equiv [\ket{g_1},\dots,\ket{g_d}] \in \mathbb{C}^{\dim_n\times d}.
\]
Then
\begin{equation}
A^{(+)}_p \equiv T^{(+)}_p\, G \in \mathbb{C}^{N_{n+1}\times d},\qquad
A^{(-)}_q \equiv T^{(-)}_q\, G \in \mathbb{C}^{N_{n-1}\times d}.
\end{equation}

\section{2nd-order contraction formula (two-site effective Hamiltonian)}
Let the intermediate energies be $E_{n+1}(\nu)$ and $E_{n-1}(\rho)$, and the two-site ground energy reference be $E_0$ (from CEF ground energies).

Define the hopping amplitude (channel $j\to i$ then back):
\begin{equation}
X_{\nu\rho}^{(c d)}=\sum_{p,q} t_{pq}\, A^{(+)}_p[\nu,c]\; A^{(-)}_q[\rho,d],
\end{equation}
where $c,d\in\{1,\dots,d\}$ label the ground-manifold basis on sites $i$ and $j$ respectively.

Then the effective Hamiltonian matrix elements in the product ground basis
$\ket{g_c}_i\otimes\ket{g_d}_j$ are
\begin{equation}
\left(H_{\text{eff}}^{(2)}\right)_{(cd),(ab)}
=
\sum_{\nu,\rho}
\frac{
X_{\nu\rho}^{(cd)}\; \left(X_{\nu\rho}^{(ab)}\right)^*
}{
E_0 - \left(E_{n+1}(\nu)+E_{n-1}(\rho)\right)
}.
\end{equation}
This is the core contraction implemented numerically (typically via \texttt{einsum}).

\section{Practical caching scheme for parameter scans}
\begin{itemize}[leftmargin=2em]
\item \textbf{Cache 1 (parameter-independent):} sparse creation matrices $D_p^{(n)}$ on the determinant basis (for each $n$).
\item \textbf{Cache 2 (parameter-dependent):} projected transitions $T_p^{(+)}=V_{n+1}^\dagger D_p^{(n)}$ and $T_q^{(-)}=V_{n-1}^\dagger C_q^{(n)}$.
\item \textbf{Intermediate-state consistency:} to safely reuse projected transitions, record the absolute path(s) of the LSJM state files used to build $V_{n\pm1}$ (e.g.\ \texttt{source\_paths.txt}); reuse only if the paths and parameter keys match exactly.
\end{itemize}

\section{Summary of the computational pipeline}
\begin{enumerate}[leftmargin=2em]
\item Single-site stage: compute and write to disk
  \begin{itemize}
  \item LSJM eigenstates (Fock expansions) and energies for $f^{n-1}, f^{n+1}$ (intermediate sectors),
  \item CEF ground manifold (Fock expansions) for $f^n$.
  \end{itemize}
\item SOPT stage: read these files and compute
\[
D_p^{(n)} \Rightarrow T^{(\pm)} \Rightarrow A^{(\pm)} \Rightarrow X \Rightarrow H_{\text{eff}}^{(2)}.
\]
\end{enumerate}

\end{document}
